\section{Convertir las matemáticas a Lean: el ciclo cerrado de la demostración formal}

Convertir las matemáticas a Lean no es una simple traducción. Las demostraciones en lenguaje natural suelen omitir el contexto, saltarse pasos, apoyarse en la intuición geométrica o citar teoremas de forma imprecisa. Lean exige explicitar el tipo de cada objeto, el enunciado del teorema, las dependencias, las tactic y la validez de cada paso. Esto le da a la IA tanto oportunidades como restricciones.

\begin{figure}[H]
\centering
\includegraphics[width=0.92\textwidth]{lean-proof-loop.png}
\caption{Ciclo cerrado desde la demostración en lenguaje natural hasta la verificación en Lean.}
\end{figure}

\begin{knowledgebox}{Lectura de la figura: por qué el ciclo cerrado de Lean es adecuado para el entrenamiento}
La demostración en lenguaje natural se convierte primero en un formal statement; después se generan las tactic o el proof term, y el kernel de Lean comprueba si son válidos. Si fallan, el sistema devuelve el error/goal state, y una persona o el modelo hace la corrección a partir de ello. Este ciclo cerrado proporciona una retroalimentación de entrenamiento más sólida que la generación de texto ordinaria.
\end{knowledgebox}

\subsection{Tres tipos de dificultades de la formalización}

La primera es la dificultad semántica: expresiones del lenguaje natural como «es evidente» o «de aquí se sigue» deben desarrollarse por completo. La segunda es la dificultad de las dependencias de la biblioteca: en Lean, los nombres y la forma de uso de los teoremas, las definiciones y los lemma deben ser exactos. La tercera es la dificultad de búsqueda: el espacio de demostraciones es enorme, y el modelo debe decidir qué subobjetivo demostrar primero, qué lemma invocar y cuándo reescribir el objetivo.

\begin{warningbox}{Lean no hace que el modelo entienda matemáticas por sí solo}
Lean puede verificar si una demostración es válida, pero no elige por el modelo los problemas importantes ni le proporciona por sí solo un sentido estético matemático. AI for Math también necesita la intuición, la selección de problemas y el criterio teórico de los matemáticos humanos.
\end{warningbox}

\subsection{Las demostraciones del Libro celestial de las matemáticas}

El título de la entrevista menciona «las demostraciones del Libro celestial de las matemáticas», que pueden entenderse como aquellas demostraciones elegantes y breves que parecen salidas de un libro celestial. Si la IA solo hace búsqueda por fuerza bruta, puede producir demostraciones largas; solo si aprende la intuición y el sentido estético matemáticos podrá acercarse a este tipo de demostraciones. La verificación formal garantiza la corrección, pero la elegancia sigue requiriendo un juicio estético.

\begin{importantbox}{Corrección y elegancia}
La verificación por máquina resuelve si una demostración «es correcta», pero no resuelve directamente si «es buena». La investigación matemática también se interesa en si la demostración explica el fenómeno, si revela una estructura y si es generalizable. El objetivo a largo plazo de AI for Math debe buscar a la vez la corrección, la legibilidad y la comprensión profunda.
\end{importantbox}

\subsection{Resumen de la sección}

Lean convierte las demostraciones en objetos ejecutables y verificables. Lo esencial de AI for Math es conectar la intuición del lenguaje natural con el ciclo cerrado de verificación por máquina, sin perder el sentido estético matemático.
